\documentclass[11pt,a4paper]{article}

\usepackage[T1]{fontenc}
\usepackage[utf8]{inputenc}
\usepackage{lmodern}
\usepackage{microtype}
\usepackage[a4paper,margin=27mm,headheight=14pt]{geometry}
\usepackage{amsmath,amssymb,amsthm,mathtools}
\usepackage{xcolor}
\usepackage{fancyhdr}
\usepackage[colorlinks=true,linkcolor=blue!45!black,
  citecolor=blue!45!black,urlcolor=blue!50!black]{hyperref}

\definecolor{paperblue}{RGB}{25,65,101}
\definecolor{softblue}{RGB}{241,247,252}
\hypersetup{
  pdftitle={One Map, Three Descriptions},
  pdfauthor={Bartosz Naskrecki, repository compilation and independent audit},
  pdfsubject={Marked roots of binary cubics and a three-dimensional Keller map},
  pdfkeywords={Jacobian conjecture, Keller map, binary cubic, marked root, resultant}
}

\pagestyle{fancy}
\fancyhf{}
\fancyhead[L]{\small\color{paperblue}One map, three descriptions}
\fancyhead[R]{\small\color{paperblue}Marked roots of binary cubics}
\fancyfoot[C]{\thepage}
\renewcommand{\headrulewidth}{0.35pt}

\newtheorem{theorem}{Theorem}[section]
\newtheorem{proposition}[theorem]{Proposition}
\newtheorem{lemma}[theorem]{Lemma}
\theoremstyle{remark}
\newtheorem{remark}[theorem]{Remark}

\DeclareMathOperator{\Sym}{Sym}
\DeclareMathOperator{\Res}{Res}
\newcommand{\A}{\mathbb A}
\newcommand{\C}{\mathbb C}
\newcommand{\Gm}{\mathbb G_{\mathrm m}}
\newcommand{\PP}{\mathbb P}
\newcommand{\dd}{\mathrm d}

\newenvironment{summarybox}
  {\par\medskip\noindent\begin{minipage}{\textwidth}%
   \color{paperblue}\hrule\color{black}\vspace{0.65em}}
  {\vspace{0.35em}\color{paperblue}\hrule\color{black}%
   \end{minipage}\par\medskip}

\title{\color{paperblue}\bfseries
  One Map, Three Descriptions\\[2mm]
  \large Marked roots of binary cubics and a three-dimensional Keller map}
\author{Repository compilation, exposition, and independent audit\\
  Bartosz Naskrecki}
\date{20 July 2026}

\begin{document}

\maketitle

\begin{abstract}
An explicit polynomial map $F\colon\C^3\to\C^3$ has constant Jacobian
determinant $-2$ but sends three distinct points to one point.  The formula
looks mysterious; its geometry is not.  A general binary cubic has three
distinct projective roots.  Mark one simple root and then forget which root
was marked.  A special
affine slice makes both the space of cubics and the space of marked cubics
isomorphic to affine three-space.  This note unifies three public
presentations of that map: the formula posted by Levent Alp\"oge, with credit
to Fable; the symmetric-product description posted by Andy Jiang
(@davikrehalt), explicitly presented there as GPT output; and Aaron Lou's
factorization-resultant derivation.  We give the exact coordinate changes,
a short proof of etaleness, the complete set-theoretic fiber counts, and the
mechanism by which inverse branches escape to infinity.
\end{abstract}

\begin{summarybox}
\textbf{Exact claim.}
Define $F=(F_1,F_2,F_3)$ by
\begin{align*}
F_1&=(1+xy)^3z+y^2(1+xy)(4+3xy),\\
F_2&=y+3x(1+xy)^2z+3xy^2(4+3xy),\\
F_3&=2x-3x^2y-x^3z.
\end{align*}
Then
\[
  \det JF=-2
\]
and
\[
F\left(0,0,-\frac14\right)
=F\left(1,-\frac32,\frac{13}{2}\right)
=F\left(-1,\frac32,\frac{13}{2}\right)
=\left(-\frac14,0,0\right).
\]
Thus $F$ is locally invertible everywhere but is not globally injective.
These are exact polynomial identities, not numerical approximations.
\end{summarybox}

The Jacobian conjecture says that a polynomial self-map of $\C^n$ with
nonzero constant Jacobian determinant must have a polynomial inverse.  By the
standard complex Jacobian conjecture, the two identities in the box therefore
constitute a counterexample in dimension three.  They are independently
reproducible, but at the date of this note the announcement had not yet passed
through the usual refereeing process.

\section{Provenance and the three viewpoints}

The public record available on 20 July 2026 separates naturally into the
following roles.

\begin{itemize}
\item \textbf{Alp\"oge--Fable: the coordinate map.}
  Levent Alp\"oge's post at 02:19:17 UTC gives the formula, determinant, and
  three-point fiber \cite{Alpoge}.  The post credits Akhil for asking about
  the problem and Fable for working during the World Cup final.  We use the
  label ``Alp\"oge--Fable presentation'' without trying to infer a more
  detailed private division of authorship.
\item \textbf{Jiang--GPT: the projective geometry.}
  At 12:02:19 UTC, Andy Jiang (@davikrehalt) posted the symmetric-product
  construction below and explicitly introduced it in the post as GPT output
  \cite{Jiang}.
\item \textbf{Lou: the resultant chart.}
  Aaron Lou gave a self-contained derivation by factoring a cubic into a
  linear and a quadratic factor, normalizing their resultant, and finding a
  global affine chart \cite{Lou}.
\end{itemize}

This repository compiles, compares, and independently checks those public
descriptions.  Its maintainer does not claim discovery authorship for any of
them.  The three descriptions fit into one diagram:

\[
\begin{array}{ccccc}
\A^3
&\xrightarrow{\ \sim\ }&
M_{\omega,\tau}
&\xrightarrow{\ \sim\ }&
X_\tau\simeq X_{\mathrm{proj}}
\\[-1mm]
{\scriptstyle F\ \text{or}\ G}\downarrow
&&\downarrow{\scriptstyle (\ell,q)\mapsto\ell q}
&&\downarrow{\scriptstyle\text{forget the marked root}}
\\[-1mm]
\A^3
&\xrightarrow{\ \sim\ }&
H_\tau
&\xrightarrow{\ \sim\ }&
Y_{\mathrm{proj}}.
\end{array}
\tag{1.1}
\]

From left to right, these are the explicit formula, Lou's normalized factor
space, the affine marked-root incidence, and Jiang's projective open set.
The horizontal arrows are isomorphisms, but they require choices of bases,
signs, and affine levels.  Thus they are isomorphic presentations of the same
abstract algebraic morphism; they are not literally the same coordinate
formulas.

For a nonexpert, four translations are useful.  A \emph{binary cubic} is a
homogeneous degree-three polynomial in two variables.  Its roots are points
of the projective line $\PP^1$.  A root is \emph{simple} when it has
multiplicity one.  An \emph{etale} map has no local folding: its derivative is
invertible everywhere.  A \emph{nonproper} map may nevertheless lose an
inverse branch because source points can run to infinity while their images
remain bounded.

\section{Jiang--GPT: mark a root, then forget it}

An effective divisor of degree $d$ on $\PP^1$ is an unordered list of $d$
points, counted with multiplicity.  Equivalently, it is the zero set of a
binary form of degree $d$, up to nonzero scale.  Hence
\[
  \Sym^d(\PP^1)\simeq\PP^d.
\]
Consider the projective morphism
\[
\mu\colon\PP^1\times\Sym^2(\PP^1)\longrightarrow\Sym^3(\PP^1),
\qquad (p,\{q,r\})\longmapsto\{p,q,r\}.
\tag{2.1}
\]
It remembers a cubic divisor but forgets which point was distinguished.  If
$p,q,r$ are distinct, its fiber consists of the three possible markings
\[
  (p,\{q,r\}),\qquad(q,\{p,r\}),\qquad(r,\{p,q\}).
\]
This is the conceptual source of generic degree three and noninjectivity.

Let $R$ be the ramification divisor of $\mu$.  It consists exactly of pairs
for which the marked point also occurs in the residual quadratic:
\[
  R=\{(p,\{q,r\}):p=q\text{ or }p=r\}.
\tag{2.2}
\]
The condition $q=r$ by itself causes no ramification when $p$ is different.

The small diagonal $p\mapsto3p$ is a twisted cubic in
$\Sym^3(\PP^1)\simeq\PP^3$.  Choose a hyperplane $H_\infty$ tangent but not
osculating to that curve: its intersection divisor with the curve has type
$2p+q$, where $p\ne q$.  Set
\begin{align*}
Y_{\mathrm{proj}}&=\Sym^3(\PP^1)\setminus H_\infty,\\
X_{\mathrm{proj}}&=
\bigl(\PP^1\times\Sym^2(\PP^1)\bigr)
\setminus\bigl(R\cup\mu^{-1}(H_\infty)\bigr).
\end{align*}
Since the complement of a hyperplane in $\PP^3$ is affine space,
$Y_{\mathrm{proj}}\simeq\A^3$.  The surprising statement is that this
particular $X_{\mathrm{proj}}$ is also $\A^3$.

Here is an intrinsic affine version.  Let $W$ be two-dimensional and choose a
tensor $\tau\in\Sym^3W$ of multiplicity type $(2,1)$.  The projective
hyperplane $H_\infty$ is defined by
$\langle\tau,C\rangle=0$.  Its complement is identified with the affine slice
\[
H_\tau=\{C\in\Sym^3W^\vee:\langle\tau,C\rangle=1\}\simeq\A^3,
\qquad [C]\longmapsto\frac{C}{\langle\tau,C\rangle}.
\tag{2.3}
\]
Define the simple-root incidence
\[
X_\tau=\{(C,[r])\in H_\tau\times\PP(W):C(r)=0, \dd C_r\ne0\}.
\tag{2.4}
\]
Deleting $R$ says precisely that the marked root is simple, so (2.3) gives
natural isomorphisms $X_{\mathrm{proj}}\simeq X_\tau$ and
$Y_{\mathrm{proj}}\simeq H_\tau$, under which $\mu$ becomes
$(C,[r])\mapsto C$.

\section{Lou: fix the factorization gauge by a resultant}

Fix a nonzero alternating form $\omega\in\bigwedge^2W^\vee$.  A marked root
of $C$ is the same as a factorization
\[
  C=\ell q,\qquad \ell\in W^\vee,\qquad q\in\Sym^2W^\vee,
\]
where $\ell$ vanishes at the marked point.  This factorization has a scaling
ambiguity
\[
  (\ell,q)\longmapsto(s\ell,s^{-1}q),\qquad s\in\Gm.
\tag{3.1}
\]
Let $v_\ell\in W$ be defined by $\iota_{v_\ell}\omega=\ell$, and put
\[
  \Res_\omega(\ell,q)=q(v_\ell).
\tag{3.2}
\]
This is the linear-quadratic resultant.  It is nonzero exactly when the
marked root is simple, and it has weight one under (3.1):
\[
\Res_\omega(s\ell,s^{-1}q)=s\Res_\omega(\ell,q).
\]
Therefore every coprime scaling orbit has a unique representative with
resultant one.  Define
\[
M_{\omega,\tau}=
\{(\ell,q):\Res_\omega(\ell,q)=1,
\ \langle\tau,\ell q\rangle=1\}.
\tag{3.3}
\]
Multiplication gives an isomorphism
\[
M_{\omega,\tau}\xrightarrow{\sim}X_\tau,
\qquad(\ell,q)\longmapsto(\ell q,\PP(\ker\ell)).
\tag{3.4}
\]

The same normalization has a useful root-gradient form.  If
$r=v_\ell$, then
\[
  \dd(\ell q)_r=q(r)\ell
  =\Res_\omega(\ell,q)\,\iota_r\omega.
\tag{3.5}
\]
Thus resultant one means $\dd C_r=\iota_r\omega$.  This single equation says
both that $r$ is a root, by Euler's identity, and that it is simple.

\begin{proposition}[Why there is no local folding]\label{prop:etale}
The forgetful map $X_\tau\to H_\tau$ is etale.
\end{proposition}

\begin{proof}
Consider
\[
\Psi(\ell,q)=\bigl(\ell q,\Res_\omega(\ell,q)\bigr).
\]
If $\ell$ and $q$ are coprime and
$\dot\ell q+\ell\dot q=0$, unique factorization gives
$(\dot\ell,\dot q)=a(\ell,-q)$.  This is exactly the infinitesimal direction
of the gauge (3.1).  Along that direction the derivative of the resultant is
$a\Res_\omega(\ell,q)$, which is nonzero unless $a=0$.  Hence $\dd\Psi$ is an
isomorphism.  The map in the proposition is its base change to the two levels
in (3.3), so it is etale.
\end{proof}

It remains to see the exceptional global affine chart.  Choose an adapted
basis $T,S$ and write
\[
\ell=xT+\beta S,\qquad
q=\gamma T^2+\delta TS+\epsilon S^2.
\]
The two equations in (3.3) become
\[
x\delta+\beta\gamma=1,\qquad
x^2\epsilon-x\beta\delta+\beta^2\gamma=1.
\tag{3.6}
\]
The following formulas define an isomorphism $\A^3_{x,y,z}\to M_{\omega,\tau}$:
\begin{align*}
\beta&=1+xy,&
\gamma&=\frac{2-3xy+x^2z}{2},\\
\delta&=\frac{y+3xy^2-xz-x^2yz}{2},&
\epsilon&=-z(1+xy)^2+y^2(4+3xy).
\tag{3.7}
\end{align*}
Its inverse is polynomial and is evaluated in this order:
\[
  x=x,\qquad y_*=2\beta\delta-x\epsilon,\qquad
  z_*=4y_*\delta+2\gamma y_*^2-\epsilon.
\tag{3.8}
\]
Substitution verifies (3.6).  Conversely, (3.6) first gives
$\beta=1+xy_*$ and then recovers the remaining three formulas in (3.7).
Thus both source and target of the forgetful map are globally affine
three-space; no divisor was silently removed from the chart.

\section{Alp\"oge--Fable: recovering the displayed formula}

The direct projective derivation uses a point $p=[u:v]$, its linear factor,
and a residual quadratic
\[
L=vT-uS,\qquad Q=AT^2+BTS+CS^2.
\]
The two projective scalings have unique representatives satisfying
\[
  h:=vB-uA=-2,\qquad r:=Q(u,v)=2.
\tag{4.1}
\]
This is exactly Lou's normalization in different signs and units.  Indeed set
\[
  \ell(T,S)=L(T,-S),\qquad q(T,S)=\frac12Q(T,-S).
\tag{4.2}
\]
Then
\[
x=v,\quad\beta=u,\quad\gamma=\frac A2,\quad
\delta=-\frac B2,\quad\epsilon=\frac C2,
\tag{4.3}
\]
and the two left sides of (3.6) are respectively $-h/2$ and $r/2$.

In the direct chart one writes
\begin{align*}
u&=1+xy,&v&=x,\\
A&=2-3xy-x^2z,\\
B&=-y-xz-3xy^2-x^2yz,\\
C&=2\bigl((1+xy)^2z+y^2(4+3xy)\bigr).
\tag{4.4}
\end{align*}
Equations (4.1) hold identically.  Comparing (4.4) with (3.7) shows that
$z_{\mathrm{Lou}}=-z_{\mathrm{direct}}$.  Moreover,
\[
LQ=cT^3-2T^2S+bTS^2-2aS^3,
\qquad
a=\frac{uC}{2},\quad b=vC-uB,\quad c=vA.
\tag{4.5}
\]
Substitution in (4.5) is precisely the map $F$ displayed on the first page.

If $G$ denotes Lou's target order, define
\[
S_0(x,y,z)=(x,y,-z),\qquad R_0(a,b,c)=(c,b,a).
\]
Then the exact relation is
\[
  \boxed{G=R_0\circ F\circ S_0.}
\tag{4.6}
\]
This is left-right linear equivalence: source and target coordinates change
independently.  In particular, the displayed identification is not a
dynamical conjugacy by one common coordinate change.

Proposition~\ref{prop:etale} shows that $\det JF$ never vanishes.  A
nowhere-zero polynomial on $\C^3$ is constant, and the linear part of $F$ at
the origin is
\[
  (x,y,z)\longmapsto(z,y,2x).
\]
Its determinant is $-2$, proving $\det JF=-2$ without a large expansion.
The three substitutions on the first page then prove noninjectivity.

There is also a compact certificate that the formula really marks a root.
For $F(x,y,z)=(a,b,c)$ set
\[
Q_{a,b,c}(U,V)=cU^3-2U^2V+bUV^2-2aV^3,
\qquad (U,V)=(1+xy,x).
\]
Direct substitution gives
\[
Q_{a,b,c}(U,V)=0,\qquad
(Q_U,Q_V)(U,V)=(2V,-2U).
\tag{4.7}
\]
The marked projective point is therefore a simple root, normalized by its
gradient.  After the sign and scale changes in (4.2)--(4.6), equation (4.7)
is the coordinate form of (3.5).

\section{Fibers, the discriminant, and escape to infinity}

Identify a target point with its binary cubic as in (4.7).  Its fiber is then
read directly from the cubic's roots:
\[
\#F^{-1}(C)=
\begin{cases}
3,&C\text{ has three simple roots},\\
1,&C\text{ has a double root and a distinct simple root},\\
0,&C\text{ has a triple root}.
\end{cases}
\tag{5.1}
\]
In the middle case only the simple root can be marked; at a triple root every
possible marking lies in the deleted ramification divisor.

Before deleting ramification, the full incidence
\[
\overline X_\tau=\{(C,[r])\in H_\tau\times\PP(W):C(r)=0\}
\longrightarrow H_\tau
\]
is finite flat of degree three: every nonzero binary cubic has a projective
zero divisor of length three, counted with multiplicity.  Its branch locus is
the binary-cubic discriminant.  In the original target coordinates a
convenient normalization is
\[
\Delta(a,b,c)=b^2-cb^3-16a+18abc-27a^2c^2.
\tag{5.2}
\]
Off $V(\Delta)$, no point was deleted, so $F$ is finite and therefore proper
there.  Every discriminant cubic has a repeated-root marking in $R$, and that
boundary point is a limit of simple markings.  A sequence of simple markings
converging to it cannot remain bounded in the global $\A^3$ chart: otherwise
it would have an affine limit, although the limiting marking was deleted.
Thus an inverse branch escapes to infinity over every point of $V(\Delta)$,
whereas no branch can do so off it.  Hence the nonproperness set is exactly
$V(\Delta)$.

The image is larger than the complement of the discriminant.  Double-root
cubics retain their one simple-root preimage.  Only triple-root cubics are
missing.  In $(a,b,c)$ coordinates this omitted curve is
\[
  3bc=4,\qquad27ac^2=4,
\]
or, equivalently,
\[
  (a,b,c)=\left(\frac{t^2}{3},2t,\frac{2}{3t}\right),
  \qquad t\in\C^*.
\tag{5.3}
\]
This distinction is the whole mechanism: the map has no finite critical
point, but some inverse branches disappear through the boundary.  It is
etale and quasi-finite, yet neither finite, proper, surjective, nor injective.

\section{Scope, history, and reproducibility}

Keller's 1939 paper is the traditional source of the conjecture
\cite{Keller}; Bass, Connell, and Wright established the standard reduction
framework \cite{BCW}, and van den Essen's monograph is a broad reference
\cite{Essen}.  Jelonek developed the modern algebraic study of nonproperness
sets of polynomial maps \cite{Jelonek}.  For historical context immediately
before the announcement, see Rodr\'iguez D\'iaz \cite{Rodriguez}; contemporary
literature still described the conjecture as open even in dimension two
\cite{Bisi}.

The construction here concerns $\C^3$.  Adjoining identity coordinates gives
$F\times\mathrm{id}_{\C^{n-3}}$ in every dimension $n\ge3$.  Nothing in this
argument decides the separate two-dimensional case.  Visualizations in the
repository show the real slice $\mathbb R^3\subset\C^3$; the full complex
space has six real dimensions, so a three-dimensional animation is an
illustration rather than the whole geometry.

Exact symbolic certificates, source code, and this paper are maintained at
\url{https://github.com/nasqret/jacobian-counterexample}.  The mathematical
text and original repository graphics are released under the Creative
Commons Attribution 4.0 International license.  Third-party posts and papers
remain under their respective terms and are linked rather than republished.

\begin{thebibliography}{9}

\bibitem{Alpoge}
L.~Alp\"oge,
\newblock X post announcing the explicit polynomial map, determinant, and
  three-point fiber, crediting Akhil for the question and Fable for the work,
\newblock 20 July 2026, 02:19:17 UTC,
\newblock \url{https://x.com/__alpoge__/status/2079028340955197566}
  (accessed 20 July 2026).

\bibitem{Jiang}
A.~Jiang (@davikrehalt),
\newblock X post giving the projective symmetric-product construction and
  explicitly presenting it as GPT output,
\newblock 20 July 2026, 12:02:19 UTC,
\newblock \url{https://x.com/davikrehalt/status/2079175065695035442}
  (accessed 20 July 2026).

\bibitem{Lou}
A.~Lou,
\newblock \emph{Deriving an Explicit Polynomial Counterexample to the
  Jacobian Conjecture: A Reproducible Cubic-Factor and Resultant Construction},
\newblock research note, 20 July 2026,
\newblock \url{https://aaronlou.com/jacobian_counterexample_derivation.pdf}.

\bibitem{Keller}
O.-H.~Keller,
\newblock Ganze Cremona-Transformationen,
\newblock \emph{Monatsh. Math. Phys.} \textbf{47} (1939), 299--306,
\newblock \url{https://doi.org/10.1007/BF01695502}.

\bibitem{BCW}
H.~Bass, E.~H. Connell, and D.~Wright,
\newblock The Jacobian conjecture: reduction of degree and formal expansion
  of the inverse,
\newblock \emph{Bull. Amer. Math. Soc. (N.S.)} \textbf{7} (1982), 287--330,
\newblock \url{https://doi.org/10.1090/S0273-0979-1982-15032-7}.

\bibitem{Essen}
A.~van den Essen,
\newblock \emph{Polynomial Automorphisms and the Jacobian Conjecture},
\newblock Progress in Mathematics 190, Birkh\"auser, 2000.

\bibitem{Jelonek}
Z.~Jelonek,
\newblock The set of points at which a polynomial map is not proper,
\newblock \emph{Ann. Polon. Math.} \textbf{58} (1993), 259--266,
\newblock \url{https://eudml.org/doc/262325}.

\bibitem{Rodriguez}
L.~O. Rodr\'iguez D\'iaz,
\newblock On the origin of the Jacobian conjecture,
\newblock \emph{Comptes Rendus Math\'ematique} \textbf{364} (2026), 363--370,
\newblock \url{https://doi.org/10.5802/crmath.831}.

\bibitem{Bisi}
E.~Bisi, P.~Dyszewski, N.~Gantert, S.~G.~G. Johnston, J.~Prochno,
and D.~Schmid,
\newblock Random planar trees and the Jacobian conjecture,
\newblock \emph{J. London Math. Soc.} \textbf{113} (2026), e70416,
\newblock \url{https://doi.org/10.1112/jlms.70416}.

\end{thebibliography}

\end{document}
